% Extensión estimada: 4--4.5 páginas; cuatro subsecciones y dos tablas.
% Fuentes: HAM10000 (metadatos, Codigo/data/raw); y_true del test en Codigo/outputs/results/swin_tiny_sweep_t06_les_cwnone_predictions.json;
% Codigo/docs/ideas/criterio_seleccion_plan_20260917.md; Codigo/docs/resul/split_lesion_resul_20260926.md (§4: con el
% criterio anterior, ConvNeXt t06 tenía sens. macro 0.700 frente a 0.691 de t02);
% Codigo/docs/resul/revision_experimentacion_20261003.md (§2.3, §2.8 y §5); src/evaluation/clinical_ner.py;
% scripts/melanoma_detection.py; scripts/lesion_split_report.py; scripts/statistical_tests.py (bootstrap_ci,
% paired_bootstrap_diff, mcnemar_exact); scripts/ensemble_bootstrap.py (paired_bootstrap_diff_flags).
% Revisión del autor del 18/10/2026: sin fechas ni historia del proceso; cada métrica con su justificación; las
% pruebas estadísticas vienen de la antigua Sección 3.9; la especificidad de 0.85 se justifica como decisión de diseño.

\section{Criterios de evaluación}
\label{sec:met_protocolo}

Fijados los datos y la partición, quedan las reglas con que se evalúa cada modelo: qué métricas resumen su desempeño,
contra qué referencias se leen y con qué pruebas se sostienen las diferencias. La clasificación y la generación de texto
no se miden igual, así que cada tarea justifica sus métricas por separado.

\subsection{Métricas de clasificación}
\label{sec:met_seleccion}

Las métricas de clasificación tienen que reflejar que las confusiones no cuestan lo mismo. Confundir \gls{mel} con
\gls{nv} retrasa la extirpación de una lesión cuya letalidad depende del estadio en que se diagnostica, mientras que
confundir \gls{nv} con \gls{mel} lleva a extirpar sin necesidad una lesión benigna, y las confusiones entre clases
benignas tienen menos consecuencias~\cite{tschandl2018ham10000,nci2026melanoma}. Como además \gls{nv} reúne dos tercios
de \emph{HAM10000} (Sección~\ref{sec:met_ham}), un buen resultado global puede esconder un mal desempeño en las clases
que más importan. Por lo tanto, las métricas tienen que evaluar cada clase por separado y darles a todas el mismo peso.

Con esa condición, el trabajo usa dos métricas que promedian las siete clases sin ponderarlas
(Sección~\ref{sec:marco_eval_clasif}). La sensibilidad macro (Ecuación~\ref{eq:marco_acc_bal}) es el resultado
principal, porque mide qué fracción de cada clase reconoce el modelo y es la métrica principal de los trabajos
revisados en la Sección~\ref{sec:estado_arte_diagnostico}, con los que así se puede comparar. El F1 macro
(Ecuación~\ref{eq:marco_f1}), en cambio, es el criterio con que se elige la configuración de cada modelo, porque
corrige lo que la sensibilidad macro no ve: un modelo que asigna muchas imágenes a las clases minoritarias sube su
sensibilidad aunque se equivoque en las demás, y el F1 macro lo penaliza a través de la precisión de cada clase. Las
dos se informan juntas, y cuando divergen, la diferencia muestra cuánto gana el modelo en detectar las clases a costa
de dar falsas alarmas.

Para leer estas métricas se usan tres referencias. En primer lugar, el azar fija el piso de cualquier clasificador de
siete clases, una sensibilidad macro de 1/7. En segundo lugar, el clasificador trivial que responde siempre NV muestra
lo que se obtiene del desbalance sin mirar la imagen: en la prueba alcanza una exactitud de 0.669 y un F1 ponderado de
0.537, pero una sensibilidad macro de 0.143 y un F1 macro de 0.115. Por eso la exactitud y el F1 ponderado, que
favorecen a NV, se informan pero no son el resultado principal. Finalmente, los clasificadores supervisados de la
Sección~\ref{sec:met_supervisados} son la línea base (Sección~\ref{sec:marco_swin}) de los \glspl{vlm}, y la primera
pregunta de investigación se responde comparándolos.

La elección de un modelo se hace en validación, en dos niveles, y la prueba no interviene en ninguno. Entre las
configuraciones de un mismo modelo se elige la de mayor F1 macro, porque en esa comparación tienen que pesar las clases
minoritarias. Dentro de cada entrenamiento, en cambio, la detención temprana y la época que se conserva se deciden por
el F1 ponderado. La razón es el tamaño de la validación: con 14 imágenes de DF y 18 de VASC
(Tabla~\ref{tab:met_tamanos}), un solo acierto más en DF sube su sensibilidad en 0.07, de modo que el F1 macro cambia
de una época a otra por pocas imágenes y podría detener el entrenamiento por azar. El F1 ponderado, dominado por las
clases con más imágenes, varía menos. Su costo es que la época conservada puede no ser la mejor para las clases
minoritarias, y por eso las prioridades clínicas se aplican en las demás decisiones: la configuración se elige por F1
macro, el resultado principal es la sensibilidad macro y el melanoma tiene su propio umbral.

Una misma configuración tampoco da siempre el mismo resultado, porque las métricas de las clases minoritarias dependen
de pocas imágenes. Para estimar esa variación, cada modelo elegido se entrena con tres semillas (42, 43 y 44), y su
resultado se informa como media $\pm$ desvío estándar muestral. Con tres semillas, sin embargo, el desvío es una
estimación gruesa, así que las comparaciones entre modelos se sostienen en las pruebas estadísticas de la
Sección~\ref{sec:met_estadistica} y no en el desvío.

La detección de melanoma (T2) necesita un criterio propio, porque si se toma como melanoma la clase de mayor
probabilidad, el costo de perder un melanoma no interviene en la decisión. Por eso el umbral sobre la probabilidad de
melanoma se fija en validación como el de mayor sensibilidad entre los que mantienen una especificidad de al menos
0.850, y se aplica sin cambios a la prueba. Así se prioriza no perder melanomas, con un tope a las falsas alarmas: en
validación, a lo sumo el 15\,\% de los casos sin melanoma se señala como sospechoso. El \gls{auc}
(Ecuación~\ref{eq:marco_auc}), que no depende del umbral, se informa como resumen de la discriminación.

\subsection{Evaluación de la generación}
\label{sec:met_eval_generacion}

A diferencia de una clase, un reporte no tiene una única respuesta correcta, y por eso el texto generado se evalúa en
dos niveles, según haya o no un texto de referencia. En los dos se evalúa solo el reporte: los bloques de razonamiento
que el generador escribe entre etiquetas \texttt{\textless{}think\textgreater{}} se descartan antes de medir, incluido
el que queda abierto cuando el límite de tokens corta el texto. Se descartan porque el clínico lee el reporte y no el
razonamiento, y porque, si se midiera el texto completo, los diagnósticos que el razonamiento solo menciona como
hipótesis contarían como aciertos.

El Nivel~1 compara el reporte con un texto de referencia mediante BLEU-4, ROUGE-L, METEOR, CIDEr y BERTScore
(Sección~\ref{sec:marco_eval_texto}). Solo puede aplicarse a \emph{MM-Skin} y \emph{SkinCoT}, los conjuntos que tienen
ese texto, y se mide sobre sus particiones de prueba, de 1\,650 y 457 imágenes. De ellas se descartan los pares sin
referencia, es decir, las imágenes de \emph{MM-Skin} sin descripción y los textos de \emph{SkinCoT} que están enteros
dentro de bloques de razonamiento, de modo que quedan 1\,414 y 426 pares. Además, cuando el límite de tokens llega
mientras el generador todavía razona, el reporte queda vacío. Esos reportes no entran en las métricas, porque no hay
texto que comparar, pero su proporción se informa junto a ellas.

Los valores del Nivel~1 se informan, pero no se usan para elegir ningún modelo ni ninguna configuración, porque los
textos de referencia no son diagnósticos clínicos de la lesión. Lo que mide la comparación depende de quién escribió
esos textos (Tabla~\ref{tab:met_conjuntos}). En \emph{MM-Skin} mide el parecido con la descripción de un libro de
texto, y en \emph{SkinCoT}, el parecido con un razonamiento que un modelo de lenguaje escribió conociendo el
diagnóstico.

El Nivel~2, en cambio, no necesita referencia y se aplica a la prueba de \emph{HAM10000}, con las dos métricas sin
referencia de la Sección~\ref{sec:marco_eval_texto}. La cobertura de entidades clínicas (\emph{clinical entity
recall}) se informa como complemento, porque mide cuánto contenido clínico describe el texto y no si es pertinente. La
consistencia diagnóstica (\emph{diagnostic consistency}, Ecuación~\ref{eq:marco_fd}) es la métrica principal de
generación, porque responde directamente la segunda pregunta de investigación: si el reporte afirma el diagnóstico
correcto de la imagen que recibe.

Para calcular la consistencia, el diagnóstico se reconoce con una búsqueda estricta, que solo acepta los términos
distintivos de cada clase, como «melanoma» o «basal cell», y exige palabra completa en las abreviaturas, como «BCC».
Una búsqueda que aceptara abreviaturas o términos genéricos, en cambio, acreditaría diagnósticos que el texto no
nombra, como «mel» dentro de «melanocytic» o la queratosis benigna por la palabra «keratosis» de «actinic keratosis».

Aun así, la consistencia diagnóstica mide fidelidad diagnóstica y no validez clínica, porque, como las demás métricas
de generación, no comprueba que lo descrito esté en la imagen (Sección~\ref{sec:marco_eval_texto}). Por eso ningún
resultado se presenta como evidencia de validez clínica. Finalmente, la Tabla~\ref{tab:met_roles} resume el papel de
cada métrica en el trabajo.

\begin{table}[H]
    \caption{Papel de cada métrica en el trabajo.}
    \label{tab:met_roles}
    \begin{tabularx}{\textwidth}{@{}>{\raggedright\arraybackslash}p{4.6cm}
        >{\raggedright\arraybackslash\hspace{0pt}}X >{\raggedright\arraybackslash\hspace{0pt}}p{3.6cm}@{}}
        \toprule
        \textbf{Métrica} & \textbf{Papel en el trabajo} & \textbf{Se define en} \\
        \midrule
        Sensibilidad macro & Resultado principal de clasificación en la prueba & Sección~\ref{sec:marco_eval_clasif} \\
        \filasep
        F1 macro & Elige la configuración de cada modelo (validación); se informa en la prueba & Ecuación~\ref{eq:marco_f1} \\
        \filasep
        F1 ponderado & Elige la época dentro de cada entrenamiento (validación) & Ecuación~\ref{eq:marco_f1} \\
        \filasep
        Exactitud & Se informa & Sección~\ref{sec:marco_eval_clasif} \\
        \filasep
        Sensibilidad y especificidad de melanoma & Umbral fijado en validación (especificidad de al menos 0.850); resultado en la prueba & Sección~\ref{sec:marco_eval_clasif} \\
        \filasep
        AUC de melanoma & Se informa & Sección~\ref{sec:marco_eval_clasif} \\
        \filasep
        Nivel 1 (BLEU-4, ROUGE-L, METEOR, CIDEr, BERTScore) & Se informa, solo en \emph{MM-Skin} y \emph{SkinCoT} & Sección~\ref{sec:marco_eval_texto} \\
        \filasep
        Cobertura de entidades clínicas & Complemento del Nivel 2 & Sección~\ref{sec:marco_eval_texto} \\
        \filasep
        Consistencia diagnóstica & Métrica principal de generación (solo el reporte, búsqueda estricta) & Sección~\ref{sec:marco_eval_texto} \\
        \filasep
        Calibración (ECE, Brier) & Se analiza en la Sección~\ref{sec:met_complementarios} & Sección~\ref{sec:marco_eval_calib} \\
        \bottomrule
    \end{tabularx}
\end{table}

\subsection{Comparación estadística}
\label{sec:met_estadistica}

Una vez calculadas las métricas, que un modelo tenga una métrica más alta que otro en la prueba no alcanza para afirmar
que es mejor (Sección~\ref{sec:marco_eval_estadistica}). Por ello, cada comparación usa el \emph{bootstrap} o la
prueba de McNemar exacta según lo que se pregunta, y la Tabla~\ref{tab:met_pruebas} indica cuál se aplica en cada caso.

Como se comparan siete clasificadores (los cuatro supervisados, DermLIP sin ajuste y los dos \glspl{vlm}
generativos sin ajuste, Sección~\ref{sec:met_clasif_vlm}), que forman 21 pares, los valores p de McNemar se corrigen
con el método de Holm (Ecuación~\ref{eq:marco_holm}). Todos los \emph{bootstrap} usan 2\,000 remuestreos de las
imágenes de prueba, con semilla 42. La sensibilidad de cada clase, en cambio, se compara de forma descriptiva, contando en cuántas de las siete
clases mejora cada modelo, porque con siete valores ninguna prueba de rangos podría alcanzar el umbral que Holm exige
para 21 comparaciones. Las pruebas se calculan sobre la ejecución con semilla 42, y la variación entre semillas se
informa como desvío. Además, el \emph{bootstrap} remuestrea imágenes y no lesiones, así que trata como independientes las imágenes de una misma
lesión; en la prueba hay 1\,503 imágenes de 1\,429 lesiones, y remuestrear por lesión casi no cambia los intervalos.

\begingroup
\small
\setlength{\tabcolsep}{3pt}
\begin{table}[H]
    \centering
    \caption{Pruebas estadísticas según la comparación.}
    \label{tab:met_pruebas}
    \begin{tabularx}{\textwidth}{@{}>{\raggedright\arraybackslash}p{3.4cm}
        >{\raggedright\arraybackslash\hspace{0pt}}X
        >{\raggedright\arraybackslash\hspace{0pt}}X
        >{\raggedright\arraybackslash\hspace{0pt}}p{2.6cm}@{}}
        \toprule
        \textbf{Comparación} & \textbf{Pregunta que responde} & \textbf{Prueba} & \textbf{Criterio} \\
        \midrule
        Métrica agregada de un modelo
            & ¿Cuánto varía por la composición de la prueba?
            & Intervalo percentil de \emph{bootstrap}
            & IC del 95\,\% \\
        \filasep
        Diferencia de F1 macro o de sensibilidad macro entre dos modelos
            & ¿La diferencia excede la variación de la muestra?
            & \emph{Bootstrap} pareado
            & IC del 95\,\% sin el cero \\
        \filasep
        Aciertos por imagen entre dos modelos
            & ¿Uno acierta imágenes en las que el otro falla?
            & McNemar exacto sobre los casos discordantes, con corrección de Holm
            & $p < 0.05$ corregido \\
        \filasep
        Sensibilidad por clase entre dos modelos
            & ¿En cuántas clases mejora cada uno?
            & Descriptiva
            & --- \\
        \filasep
        Consistencia diagnóstica entre dos variantes del \emph{pipeline}
            & ¿La diferencia excede la variación de los reportes?
            & \emph{Bootstrap} pareado de la diferencia de proporciones
            & IC del 95\,\% sin el cero \\
        \bottomrule
    \end{tabularx}
\end{table}
\endgroup

Con estos criterios se entrenan y se eligen los clasificadores supervisados de la sección siguiente.
